% le caratteristiche richieste dall'università sono elencate qui: https://stem.elearning.unipd.it/mod/book/view.php?id=234&chapterid=46#modalita
% 12pt: font richiesto dall'università
% twoside: i margini interni ed esterni sono scambiati per le pagine "a sinistra" e "a destra"
% openright: i capitoli cominciano in pagine dispari ("a destra")
% extreport: supporta 12pt

\documentclass[12pt,a4paper,twoside,openright]{extreport}

\usepackage{amsmath}                            % per avere più controllo sulle equazioni 
\usepackage{enumitem}                           % per avere più controllo sulle enumerazioni
\usepackage[
    a4paper,
    top=2cm,bottom=2cm,
    outer=2cm,inner=3cm,
    includeheadfoot
]{geometry}                                     % margini richiesti dall'università
\usepackage{graphicx}                           % per le immagini
\usepackage[newfloat]{minted}                             % per il codice con la colorazione della sintassi
\usepackage{subcaption}                         % per le sotto didascalie
\usepackage[italian]{babel}
\usepackage{csquotes}                           % per le citazioni
\usepackage[backend=biber,style=numeric,natbib=true,sorting=nyt]{biblatex}
\usepackage{fontspec}
\usepackage{setspace}
\usepackage{fancyhdr}
\usepackage{etoolbox}
\usepackage{rotating}
\usepackage[pdfa]{hyperref}
\usepackage{hyperxmp}
\usepackage{caption}


\sloppy % per evitare che il testo in \verb finisca oltre i margini
\selectlanguage{italian}
\bibliography{bibliografia}
\setmainfont{Times New Roman} % carattere richiesto dall'università
\onehalfspacing % interlinea richiesta dall'università
\setlength{\skip\footins}{2\bigskipamount} % spazio tra testo e piè di pagina
\setlength{\footnotesep}{12pt} % spazio tra ogni nota a piè di pagina

\newcommand{\bchapter}[1]{ % comando per NON mostrare l'header nelle prime pagine di ogni capitolo
  \chapter[#1]{#1\thispagestyle{empty}}
  \pagestyle{plain}
}

\renewcommand\captionfont{\footnotesize}  %imposta la grandezza delle caption come le note a piè di pagina
%\captionsetup[listing]{skip=-10pt}        %diminuisce lo spazio tra il codice e la sua caption
\AtBeginEnvironment{figure}{\vspace{1\baselineskip}} % Aggiunge spazio prima dell'ambiente figure

\AtBeginEnvironment{listing}{\vspace{1\baselineskip}} % Aggiunge spazio prima dell'ambiente listing

\newenvironment{code}{\captionsetup{type=listing}}{}  % Setto l'ambiente per l'inserimento dei codici
\SetupFloatingEnvironment{listing}{name=Codice}


% questi valori vengono usati nella composizione del frontespizio
\title{Progetto Programmazione di Sistemi Embedded: Mobile AI}
\author{Alessandro Girlanda, Andrea Mutti, Umberto Bianchin}

\hypersetup{%
  pdftitle={Progetto Programmazione di Sistemi Embedded: Mobile AI}
  pdfauthor = {Alessandro Girlanda, Andrea Mutti, Umberto Bianchin},
  pdfapart=1, %set to 1 for PDF/A-1
  pdfaconformance=B,
}

\begin{document}
    \pagenumbering{roman}
    \pagestyle{empty} % per le prime pagine, non mostrare il numero di pagina

    \begin{titlepage}
    % solo il frontespizio deve essere simmetrico rispetto ai margini interno ed esterno
    \newgeometry{hmargin=2.5cm,vmargin=2cm}
        \begin{figure}
            \centering
            \begin{subfigure}[b]{0.4\textwidth}
                \includegraphics[width=\textwidth]{Immagini/logo_unipd}
            \end{subfigure}
            \hfill
            \begin{subfigure}[b]{0.3\textwidth}
                \includegraphics[width=\textwidth]{Immagini/logo_dei.png}
            \end{subfigure}
        \end{figure}
    
        \vspace*{\stretch{0.5}}
    
        \begin{center}
            \makeatletter % serve per poter usare \@...

            % NOTA: il Times New Roman non supporta il maiuscoletto.
            %\textbf{DIPARTIMENTO DI INGEGNERIA DELL'INFORMAZIONE}\\
            %\vspace*{\stretch{0.1}}
            %\textbf{CORSO DI LAUREA IN INGEGNERIA INFORMATICA}
    
            %\vspace*{\stretch{0.5}}
            \LARGE
            \textbf{\@title}

            \begin{figure}[ht]
                \centering
                \includegraphics[width=0.5\textwidth]{Immagini/App/icona_app.jpg}
            \end{figure}
    
            \vspace*{\stretch{1.5}}
            \large
            \begin{flushright}
                \textbf{Gruppo:} \\ Alessandro Girlanda \\  Andrea Mutti \\  Umberto Bianchin \\
            \end{flushright}
    
            \vspace*{\stretch{2}}
            \textbf{ANNO ACCADEMICO 2023-2024}
        
            \makeatother % serve dopo \makeatletter
        \end{center}
    \restoregeometry
\end{titlepage}

    \cleardoublepage
    

    \pagestyle{plain} % comincia a mostrare il numero di pagina

    \tableofcontents
    \listoffigures
    %\listoflistings
    \cleardoublepage % per assicurarsi che la numerazione araba cominci col primo capitolo

    \pagenumbering{arabic}

    \fancypagestyle{plain}{                       %setto l'intestazione della pagina
    \fancyhf{} 
    \fancyhead[RO, LE]{\thepage}
    \fancyhead[LO, RE]{\nouppercase{\leftmark}}
    \renewcommand{\headrulewidth}{0.4pt}
    }
    \setlength{\headheight}{14.5pt}

    \bchapter{PyTorch Mobile}
\section{PyTorch}
Pytorch\cite{Pytorch} è un framework di deep learning open source, sviluppato inizialmente da \textit{Meta AI} e ora parte della \textit{Linux Foundation};
progettato con Python, viene usato per creare reti neurali e per progetti di apprendimento automatico,
combinando la libreria di machine learning \textbf{Torch}\cite{Torch} con un’API di alto livello basata su Python. Torch è famoso, specialmente nel campo
del Deep Learning, per fornire tool semplici e flessibili insieme a performance elevate; uno dei sui punti salienti è il grande supporto per le GPU, che 
contribuisce ad un allenamento più efficiente dei modelli di deep learning.
PyTorch fornisce un pacchetto Python per funzionalità ad alto livello, come l'elaborazione dei \textbf{tensori}\footnote{Array multidimensionale utilizzato per memorizzare dati. Nel campo del Machine Learning vengono usati per rappresentare e manipolare input, pesi e output.}
e un così detto \textbf{TorchScript}, che permette di creare modelli da PyTorch che possono poi venire salvati e caricati in un processo dove non c'è alcuna 
dipendenza di Python.

\section{PyTorch Mobile}
PyTorch Mobile si riferisce ad un set di librerie e funzionalità fornite da Pytorch. 

    \bchapter{TensorFlow}
TensorFlow è una libreria open source per l'apprendimento automatico, il calcolo numerico e altre attività di analisi statistica e predittiva.
Questo tipo di tecnologia, sviluppata e rilasciata da Google nel novembre 2015, rende l’implementazione di modelli di machine learning più semplice e
veloce per gli sviluppatori, assistendo nel processo di acquisizione dei dati, nella formulazione di previsioni su larga scala e nel successivo affinamento
dei risultati.

Lo scopo principale di TensorFlow è la creazione e l’addestramento di reti neurali, che possono essere utilizzate per moltissime applicazioni, quali:
\begin{itemize}
    \item Classificazione delle immagini;
    \item Elaborazione del linguaggio naturale;
\end{itemize}

\begin{figure}
    \centering
    \includegraphics[width=0.5\textwidth]{Immagini/logo_dei.png}
    \caption{Diagramma dei punteggi di utilizzo di vari framework nel 2018}
    \label{fig:diagramma}
\end{figure}

    
    \fancypagestyle{plain}{
    \fancyhf{}
    \renewcommand{\headrulewidth}{0pt}
    }

    \begin{spacing}{1.8}                    %spaziatura tra un elemento e l'altro della bibliografia
      \printbibliography[heading=bibintoc]
    \end{spacing}

\end{document}
